\section{پاسخ سوال ۱۰}
\begin{stepbox}
\field{تقسیم‌بندی تصویر با استفاده از Normalized Cut}{
برای اعمال برش نرمال (Normalized Cut)، ابتدا باید گراف شبکه‌ای تصویر و وزن یال‌ها را تشکیل دهیم. 
پیکسل‌های تصویر به صورت افقی و عمودی با یکدیگر ارتباط دارند. وزن هر یال با تابع گاوسی تفاوت شدت روشنایی محاسبه می‌شود: $W_{ij} = e^{-(I_i - I_j)^2 / \sigma^2}$.

با دقت در اعداد تصویر ۴×۴ داده شده، مشاهده می‌کنیم که تفاوت شدت روشنایی بین \textbf{هر دو پیکسل مجاور} (چه به صورت افقی و چه به صورت عمودی) دقیقاً برابر با \textbf{۲۰} است:
\begin{itemize}
    \item افقی: $|160 - 140| = 20, |140 - 120| = 20, \dots$
    \item عمودی: $|160 - 140| = 20, |120 - 100| = 20, \dots$
\end{itemize}

بنابراین وزن تمام یال‌های گراف شبکه‌ای یکسان و برابر با ثابت $W_{const} = e^{-400/\sigma^2}$ است. 

\textbf{نتیجه برش نرمال:}
زمانی که گراف ما یک شبکه یکنواخت \lr{(Uniform Grid)} با یال‌های کاملاً هم‌وزن باشد، روش Normalized Cut که به دنبال بهینه‌سازی فرمول $\frac{cut(A,B)}{assoc(A,V)} + \frac{cut(A,B)}{assoc(B,V)}$ است، گراف را به متقارن‌ترین حالت ممکن برش می‌دهد تا حداقلِ وزن یال‌ها قطع شود در حالی که دو حجم برابر ایجاد می‌گردد. 

بردار فیدلر \lr{(Fiedler Vector)} برای چنین گرافی، شبکه را دقیقاً \textbf{از وسط به دو نیمه مساوی} تقسیم می‌کند. این برش می‌تواند به یکی از دو صورت زیر باشد:
\begin{enumerate}
    \item \textbf{برش افقی:} جدا کردن دو سطر بالایی از دو سطر پایینی.
    \item \textbf{برش عمودی:} جدا کردن دو ستون سمت چپ از دو ستون سمت راست.
\end{enumerate}
در هر دو حالت، تعداد یال‌های قطع‌شده برابر ۴ عدد بوده و تصویر دقیقاً به دو سگمنت ساختاری یکسان ۲×۴ (یا ۴×۲) تقسیم می‌شود.
}
\end{stepbox}
